% !TEX root = main.tex
% Auto-generated by out/r5_posthoc_hparam_sensitivity_20260917/code/make_paper_sections.py
\subsection{Post-hoc Hyper-parameter Sensitivity and Cost-component Ablation}
\label{sec:posthoc-sensitivity}

\noindent\emph{As a post-hoc supplementary analysis we evaluate hyper-parameter sensitivity on the development seeds 201-205 only. This analysis was not preregistered and was not used to select any main-experiment hyper-parameter; the frozen holdout seeds 301-305 were not re-run, and the holdout markers shown in the figures are independent reference points copied read-only from pre-existing frozen results.}

\paragraph{Scope.} Three bridges (Celer cBridge, Multichain, PolyNetwork) --- identical to the confirmatory three-bridge roster --- five development seeds (201--205, 48 templates each), and two ranking methods on the \emph{same} transport plan: \texttt{RAW\_UOT\_PLAN\_D4} (mutual top-$k$ on $P$) and \texttt{CONDITIONAL\_UOT\_D4} (mutual top-$k$ on the dual-cancelled scores $S^{\mathrm{row}}_{ij}=P_{ij}/c_j$, $S^{\mathrm{col}}_{ij}=P_{ij}/r_i$). The paper defaults $k=5$, $\varepsilon=0.05$, $\lambda=0.5$ (UOT marginal relaxation $\mathrm{reg}_m$) were frozen \emph{before} this grid and are not modified by it. Grids: $k\in\{2,3,5,7,10,15\}$, $\varepsilon\in\{0.01,0.02,0.05,0.1,0.2\}$, $\lambda\in\{0.1,0.25,0.5,1,2\}$. Total: 3 bridges $\times$ 5 seeds $\times$ 14 unique configurations = 210 cells $\times$ 2 methods. The shared default configuration is solved once and referenced by all three sweeps.

\paragraph{Rank cutoff $k$.}
\begin{table}[t]
\centering\small
\caption{Development sensitivity to the rank cutoff $k$ ($\varepsilon=0.05$, $\lambda=0.5$; seeds 201--205; macro edge F1 averaged over the 15 bridge$\times$seed cells).}
\label{tab:sens-k}
\begin{tabular}{lcc}
\toprule
$k$ & \texttt{RAW\_UOT\_PLAN\_D4} & \texttt{CONDITIONAL\_UOT\_D4} \\
\midrule
2 & 0.0008 & 0.3504 \\
3 & 0.1236 & 0.4797 \\
\textbf{5} & \textbf{0.2374} & \textbf{0.3129} \\
7 & 0.2623 & 0.2820 \\
10 & 0.2594 & 0.2720 \\
15 & 0.2517 & 0.2589 \\
\bottomrule
\end{tabular}
\end{table}

The conditional decoder dominates the raw plan decoder at every $k$ (paired difference +0.3561 at $k=3$ narrowing to +0.0072 at $k=15$), so its benefit is largest exactly in the sparse regime. At $k=2$ the raw plan decoder collapses to F1 $=0.0008$ while the conditional decoder still attains $0.3504$. for \texttt{RAW\_UOT\_PLAN\_D4} the default point k=5 (F1 $=0.2374$, grid rank 4/6) is not the best point of this grid; the best point is k=7 (F1 $=0.2623$, -0.0249 relative to the default); for \texttt{CONDITIONAL\_UOT\_D4} the default point k=5 (F1 $=0.3129$, grid rank 3/6) is not the best point of this grid; the best point is k=3 (F1 $=0.4797$, -0.1668 relative to the default).

\paragraph{Entropic regularisation $\varepsilon$.}
\begin{table}[t]
\centering\small
\caption{Development sensitivity to $\varepsilon$ ($k=5$, $\lambda=0.5$; seeds 201--205) together with the measured Sinkhorn iteration counts.}
\label{tab:sens-eps}
\begin{tabular}{lcccc}
\toprule
$\varepsilon$ & RAW F1 & COND F1 & mean iters & converged \\
\midrule
0.01 & 0.2503 & 0.3136 & 630 & 15/15 \\
0.02 & 0.2455 & 0.3144 & 320 & 15/15 \\
\textbf{0.05} & \textbf{0.2374} & \textbf{0.3129} & 131 & 15/15 \\
0.1 & 0.1482 & 0.3132 & 69 & 15/15 \\
0.2 & 0.0407 & 0.3152 & 37 & 15/15 \\
\bottomrule
\end{tabular}
\end{table}

The conditional decoder is essentially flat over the whole $\varepsilon$ grid (span 0.0007), whereas the raw plan decoder degrades sharply for large $\varepsilon$ ($0.2374$ at $\varepsilon=0.05$ versus $0.0407$ at $\varepsilon=0.2$). for \texttt{RAW\_UOT\_PLAN\_D4} the default point \varepsilon=0.05 (F1 $=0.2374$, grid rank 3/5) is not the best point of this grid; the best point is \varepsilon=0.01 (F1 $=0.2503$, -0.0129 relative to the default); for \texttt{CONDITIONAL\_UOT\_D4} the default point \varepsilon=0.05 (F1 $=0.3129$, grid rank 5/5) is not the best point of this grid; the best point is \varepsilon=0.2 (F1 $=0.3152$, -0.0023 relative to the default).

\paragraph{Solver stability.}
Convergence uses the project's existing criterion (final POT marginal error $<10^{-7}$, \texttt{stopThr}$=10^{-11}$, \texttt{numItermax}$=20000$; neither was tuned). Every one of the 75 $\varepsilon$-sweep solves converged (max final error 9.98e-12) and no solve reached the iteration cap; all $k$ and $\lambda$ configurations converged as well. The iteration count falls approximately as $O(1/\varepsilon)$ (630 iterations at $\varepsilon=0.01$ versus 37 at $\varepsilon=0.2$), but the $\varepsilon$-driven degradation at large $\varepsilon$ is a property of the raw plan decoder, not of solver instability.

\paragraph{Marginal relaxation $\lambda$ and unmatched mass.}
\begin{table}[t]
\centering\small
\caption{Development sensitivity to the UOT marginal relaxation $\lambda$ ($k=5$, $\varepsilon=0.05$; seeds 201--205). $\delta^{S}_{\mathrm{total}}=\sum_i|\sum_j P_{ij}-a_i|$ and $\delta^{T}_{\mathrm{total}}=\sum_j|\sum_i P_{ij}-b_j|$ are read directly from the returned plan and the marginals actually passed to the solver.}
\label{tab:sens-lambda}
\begin{tabular}{lcccccc}
\toprule
$\lambda$ & COND F1 & RAW F1 & $\delta^{S}_{\mathrm{total}}$ & $\delta^{T}_{\mathrm{total}}$ & total mass & mean iters \\
\midrule
0.1 & 0.3130 & 0.2352 & 0.7587 & 0.7587 & 0.2413 & 33 \\
0.25 & 0.3130 & 0.2372 & 0.4839 & 0.4839 & 0.5161 & 70 \\
\textbf{0.5} & \textbf{0.3129} & \textbf{0.2374} & 0.2957 & 0.2957 & 0.7043 & 131 \\
1 & 0.3133 & 0.2376 & 0.1656 & 0.1654 & 0.8346 & 249 \\
2 & 0.3136 & 0.2377 & 0.0879 & 0.0878 & 0.9122 & 478 \\
\bottomrule
\end{tabular}
\end{table}

Raising $\lambda$ from 0.1 to 2 tightens the marginals as designed: transport mass rises from 0.2413 to 0.9122 and $\delta_{\mathrm{total}}$ falls from 1.5174 to 0.1757. Edge F1, however, is insensitive (conditional span 0.0007). On this benchmark $\lambda$ is therefore a numerically effective but metric-insensitive knob. for \texttt{RAW\_UOT\_PLAN\_D4} the default point \lambda=0.5 (F1 $=0.2374$, grid rank 3/5) is not the best point of this grid; the best point is \lambda=2 (F1 $=0.2377$, -0.0003 relative to the default); for \texttt{CONDITIONAL\_UOT\_D4} the default point \lambda=0.5 (F1 $=0.3129$, grid rank 5/5) is not the best point of this grid; the best point is \lambda=2 (F1 $=0.3136$, -0.0007 relative to the default).

\paragraph{Cost-component leave-one-out ablation.}
\begin{table}[t]
\centering\small
\caption{Cost-component leave-one-out ablation for \texttt{CONDITIONAL\_UOT\_D4} at the paper defaults (no re-tuning). Each omitted component has its weight set to zero and the remaining weights renormalised to sum to one (scale preserving, so $\varepsilon$ keeps its meaning). The baseline is the paper's primary amount-free renormalised cost, which is the ``amount omitted'' row. Seeds 201--205.}
\label{tab:cost-ablation}
\begin{tabular}{lcccc}
\toprule
Omitted & Celer $\Delta$F1 & Multichain $\Delta$F1 & PolyNetwork $\Delta$F1 & seeds $<0$ \\
\midrule
none (all six) & -0.0815 & -0.1176 & -0.0737 & 5/5/5 \\
time & -0.1830 & -0.1689 & -0.1829 & 5/5/5 \\
route & -0.0864 & -0.1267 & -0.0746 & 5/5/5 \\
risk & -0.0834 & -0.1224 & -0.0741 & 5/5/5 \\
evidence & -0.0804 & -0.1190 & -0.0722 & 5/5/5 \\
address novelty & -0.0824 & -0.1247 & -0.0720 & 5/5/5 \\
\bottomrule
\end{tabular}
\end{table}

Temporal cost is by far the most important term: removing it costs 0.1830--0.1689 macro edge F1, consistently across all five seeds and all three bridges, roughly 1.9$\times$ the effect of any other component. Route, risk, evidence and address novelty each move F1 by less than 0.01. Reinstating the amount component as a sixth term reduces F1 relative to the amount-free primary cost (-0.0815 to -0.1176), consistent with the paper's existing amount-free design.

\paragraph{Bridge specificity.} For route, risk, evidence and address novelty the ablation loss is larger on Multichain than on either other bridge (about 1.5--1.8$\times$), with the same sign in all five seeds. With only five seeds the smallest attainable two-sided paired permutation $p$-value is $0.0625$, so none of these contrasts reaches $p<0.05$, and for temporal cost Multichain is in fact the \emph{least} affected bridge. There is therefore no single Multichain-exclusive core component. The pattern \emph{is consistent with} the reading that Multichain depends a little more on the non-temporal cost terms, but it does \emph{not} establish that claim; the evidence is limited by the five-seed development sample and these are leave-one-out associations, not causal decompositions.

\paragraph{Read-only holdout reference.} The frozen confirmatory holdout (seeds 301--305) was not re-run. Its pre-existing default-configuration macro edge F1 values are used only as independent markers: Celer cBridge \texttt{CONDITIONAL} $0.3194$; Multichain \texttt{CONDITIONAL} $0.2935$; PolyNetwork \texttt{CONDITIONAL} $0.3183$. The largest absolute difference from the corresponding development default is 0.0095, but this agreement holds at the default point only and does not extend to other grid values.
